% ============================ Q7 ============================
\begin{frame}{Q7 — SGLang vs vLLM 스케줄러 · host 병목}
  \footnotesize
  \hd{How SGLang differs}
  \begin{itemize}
    \item \textbf{Overlap scheduler} — build the next batch on CPU while the current runs on device. The defining feature.
    \item Radix / cache-aware admission (longest-prefix-match first); FSM $+$ jump-forward integrated into scheduling.
    \item vLLM v1: separate EngineCore process, chunked prefill by default, token-budget scheduling.
  \end{itemize}

  \hd{Importance for RBLN — the top TPOT lever}
  \begin{itemize}
    \item Decode device step is short and fixed $\Rightarrow$ host fraction (batch build, match, block-table,
          H2D, per-step mask, sampling, detok) dominates.
    \item host $\ge$ device $\Rightarrow$ host-bound $\Rightarrow$ TPOT set by the scheduler loop, not the chip.
    \item AOT removes the graph-capture escape hatch $\Rightarrow$ overlap is close to mandatory.
    \item vllm-rbln evidence: pinned \texttt{\_to\_list}, per-step host masks, Python-loop KV copies, async scheduling disabled.
    \item Porting SGLang's overlap scheduler $=$ \textbf{greenfield} (nothing to replace).
    \item Caveat: TTFT / long prefill is device-bound — this is a decode / TPOT lever.
  \end{itemize}

  \begin{block}{Judgment basis}
    Measure device-exec vs wall step time at target batch. host critical path $>$ device exec $\Rightarrow$ scheduler is the top lever.
    \emph{[insert your own host-bottleneck example: measurement, \% host, what you changed.]}
  \end{block}
\end{frame}
